\section{Introducción}

\subsection{Contexto}

La universidad organiza cada semestre una gran cantidad de eventos abiertos a sus
estudiantes: conciertos, talleres, conferencias, ferias y actividades culturales.
Algunos se realizan en edificios o teatros del campus, cuya capacidad es conocida,
y otros en lugares externos. Un mismo evento puede repetirse en varias fechas, y
algunos son gratuitos mientras que otros tienen un costo por cupo.

Hoy la inscripción a estos eventos suele hacerse con formularios independientes,
hojas de cálculo o directamente en la entrada del evento. Esto hace difícil saber
cuántos cupos quedan, permite que se inscriban más personas que la capacidad del
lugar y no deja un registro confiable de quién reservó, modificó o canceló.

En el curso de Sistemas Distribuidos se propone construir un \textbf{sistema
distribuido de transacciones} que resuelva este problema y que, al mismo tiempo,
permita estudiar en la práctica los conceptos del curso: arquitecturas por capas,
comunicación por paso de mensajes, modelos de interacción, de fallos y de
seguridad, tolerancia a fallos, escalabilidad y evaluación de desempeño.

\subsection{Problema}

Se necesita un sistema al que muchos estudiantes accedan al mismo tiempo desde
distintas máquinas para:

\begin{itemize}
    \item consultar los eventos programados en un mes y los cupos disponibles de cada fecha;
    \item reservar uno o más cupos en una fecha de un evento;
    \item modificar una reserva (cantidad de cupos o fecha), y
    \item cancelar una reserva.
\end{itemize}

Las consultas, reservas y modificaciones deben responder de forma \textbf{síncrona}:
el estudiante espera y recibe el resultado definitivo en la misma operación. La
cancelación es \textbf{asíncrona}: el sistema la acepta de inmediato, informa que
está en proceso y la aplica después en la capa de datos.

El problema es distribuido por naturaleza: los clientes, los servicios de negocio y
el almacenamiento corren en procesos y máquinas distintas que se comunican por la
red. Esto introduce dificultades que no existen en un programa centralizado:
\begin{itemize}
    \item \textbf{Concurrencia:} dos estudiantes pueden pedir el último cupo al mismo
          tiempo y el sistema nunca debe superar la capacidad de una fecha.
    \item \textbf{Fallos parciales:} un proceso puede caerse o un mensaje perderse; un
          reintento no debe reservar ni cobrar dos veces, y una cancelación aceptada no
          debe perderse.
    \item \textbf{Desempeño:} cada transacción debe completarse en menos de un segundo,
          y el desempeño debe medirse y compararse entre una arquitectura base y una
          escalable (Tabla~1 del enunciado).
    \item \textbf{Seguridad:} se manejan datos personales de los estudiantes y dinero.
\end{itemize}

\subsection{Objetivos}

\paragraph{Objetivo general.} Diseñar, implementar y evaluar un sistema distribuido
de tres capas (consumidores, servicios de negocio y almacenamiento), comunicado con
ZeroMQ, que gestione las reservas de eventos universitarios de forma consistente,
tolerante a fallos y con tiempos de respuesta menores a un segundo.

\paragraph{Objetivos específicos.}
\begin{enumerate}
    \item Describir el sistema con los modelos arquitectónico, de interacción, de
          fallos y de seguridad, y mostrar cómo cada uno condiciona el diseño.
    \item Diseñar la arquitectura base de la Figura~\ref{fig:arquitectura-basica} con
          diagramas UML de despliegue, componentes, clases y secuencia.
    \item Implementar en C++ orientado a objetos los componentes de la arquitectura
          base con las operaciones de consulta y reserva.
    \item Identificar y justificar los componentes que harán la arquitectura escalable
          y tolerante a fallos en la segunda entrega.
    \item Definir el protocolo de pruebas del sistema, las herramientas para obtener
          las métricas de desempeño de la Tabla~1 y el método de generación de carga.
\end{enumerate}

\subsection{Alcance de esta entrega}

Este documento corresponde a la \textbf{primera entrega}. Contiene los modelos del
sistema, el diseño completo de la arquitectura base (incluidas la modificación y
la cancelación asíncrona), la justificación de los componentes de escalabilidad y
tolerancia a fallos, el protocolo de pruebas y la estrategia de medición. La
implementación entregada cubre las operaciones de \textbf{consulta y reserva}
sobre los componentes de la arquitectura base (sección~\ref{sec:implementacion}).
